% \section{Resultados}
% \label{sec:cap4_app_final_discussao}

% Os resultados apresentados neste capítulo avaliam a cadeia de aproximações da missão de \gls{rvdb}, desde a transferência orbital inicial até a aproximação final.
% Em resumo, buscou-se verificar se:
% (i) a transferência de Hohmann conduz o perseguidor a uma órbita adequada para a transição ao referencial \gls{lvlh};
% (ii) o controle translacional baseado nas equações \gls{hcw} leva o perseguidor à região de aproximação de curta distância com erro e velocidade controlados;
% (iii) o controle de atitude estabiliza o perseguidor tanto com o manipulador inoperante quanto em regime acoplado espaçonave--manipulador;
% e (iv) o manipulador posiciona a ferramenta em uma vizinhança próxima da porta de acoplamento, sem comprometer a dinâmica global do sistema.

% \subsection{Coerência entre as fases de aproximação}

% A linha do tempo da missão mostra que as três aproximações (longa, curta e final) são encadeadas de maneira consistente.
% Na aproximação de longa distância, a transferência de Hohmann eleva o raio orbital do perseguidor da altitude inicial de $500$ km para aproximadamente $599{,}9$ km, praticamente coincidente com a órbita circular do alvo, que se encontra a $600$ km de altitude. A diferença radial residual após a segunda queima é da ordem de $100$ m, valor coerente com o emprego posterior de um modelo linearizado em LVLH para descrever o movimento relativo.

% A análise da longitude inercial mostra que a condição de fase entre perseguidor e alvo é ajustada;
% ao final da transferência e do \emph{drift} orbital, a diferença de longitude reduz-se para próxmo de $0{,}0041^\circ$, o que corresponde a uma separação relativa de aproximadamente $506$ m. 

% A comparação entre os cenários (A) e (B) indica que alterações nas condições iniciais (como altitude e longitude do perseguidor) preservam a estrutura relativa da missão, mas impactam de forma sensível o tempo total e o consumo de $\Delta v$.
% No cenário B, por exemplo, o tempo até o término da aproximação de longa distância aumenta em cerca de $365\%$, e o consumo total de $\Delta v$ praticamente dobra em relação ao cenário de referência. Isso confirma que a escolha da órbita inicial e da fase inercial desempenha papel relevante no custo propulsivo da missão, ainda que o estado final para o início do movimento relativo em LVLH permaneça semelhante.

% \subsection{Desempenho do controle translacional em LVLH}

% Na aproximação de curta distância, a dinâmica relativa no referencial LVLH mostra que a lei de controle proporcional--derivativa cumpre os objetivos de redução da posição e da velocidade relativas, mantendo o perseguidor dentro de limites operacionais pré-estabelecidos. As componentes de posição partem de valores da ordem de centenas de metros e convergem para uma condição final próxima de $-0{,}007$ m em along-track e $-0{,}500$ m em radial, enquanto a componente cross-track permanece praticamente nula desde o início da fase.

% Os tempos de assentamento em posição, calculados com critério de tolerância de $\pm 0{,}05$ m durante um intervalo mínimo de $15$ s, são da ordem de $515$ s para a componente radial e $530$ s para a componente along-track, com assentamento imediato em cross-track. Esses tempos refletem um compromisso entre rapidez de convergência e limitação de esforços de controle, uma vez que a mesma lei deve respeitar restrições de velocidade relativa e de aceleração impostas para garantir segurança na vizinhança do alvo.

% A velocidade relativa inicial é dominada pela componente along-track, cerca de $0{,}1625$ m/s, e converge para valores finais da ordem de $10^{-5}$ m/s. A saturação da velocidade em $0{,}3000$ m/s para distâncias superiores a $20$ m e em $0{,}0300$ m/s na vizinhança do alvo impede que a dinâmica aproximada pelas equações de \gls{hcw} conduza a velocidades de alguns metros por segundo. Os picos de velocidade não saturada (cerca de $4{,}51$ m/s em radial e $8{,}62$ m/s em along-track) mostram que, sem essa limitação, a manobra apresentaria risco de colisão entre o perseguidor e o alvo. Em contrapartida, a aceleração de controle é mantida dentro de $\pm 1{,}0000$ m/s², e tende rapidamente a zero após o regime transitório, o que caracteriza um esforço de controle compatível com atuadores de baixa potência.
% Esses resultados indicam que o controle translacional consegue trazer o perseguidor para uma região de referência próxima ao alvo, com limites de posição e velocidade adequados para a transição à fase de aproximação final, ainda que à custa de tempos de assentamento relativamente longos.

% \subsection{Controle de atitude e dinâmica da espaçonave no modo operacional}

% Durante a missão, o controle de atitude reduz os erros de orientação do perseguidor e amortece os transitórios associados à comutação do modelo dinâmico e à atuação do manipulador robótico.
% Ao final da simulação, os ângulos de Euler convergem para valores residuais reduzidos: $0{,}10^\circ$ (roll), $0{,}02^\circ$ (pitch) e $0{,}30^\circ$ (yaw).

% As velocidades angulares da espaçonave permanecem limitadas na simulação, com faixas aproximadas de
% $[-1{,}19,\;1{,}20]$ rad/s em $w_x$, $[-1{,}26,\;1{,}25]$ rad/s em $w_y$ e $[-0{,}96,\;1{,}00]$ rad/s em $w_z$,
% convergindo ao final para $w_x=-0{,}01$ rad/s, $w_y=0{,}00$ rad/s e $w_z=-0{,}01$ rad/s.
% As acelerações angulares apresentam picos concentrados no regime transitório, com mínimo de $-109{,}94$ rad/s$^2$ em $\dot{w}_y$,
% e convergem ao final para $\dot{w}_x=-1{,}18$ rad/s$^2$, $\dot{w}_y=-2{,}00$ rad/s$^2$ e $\dot{w}_z=1{,}70$ rad/s$^2$.

% Os torques de controle de atitude permanecem tipicamente na faixa de poucos Nm, com picos absolutos próximos de $12{,}49$ Nm.
% No instante final, o torque PD de atitude assume $0{,}21$ Nm (eixo $x$), $-0{,}04$ Nm (eixo $y$) e $0{,}06$ Nm (eixo $z$).
% Na forma compensada (incluindo termos de acoplamento), os valores finais são $0{,}16$ Nm, $-0{,}51$ Nm e $0{,}08$ Nm.
% Os torques de reação na base, induzidos pela movimentação do manipulador, apresentam amplitudes inferiores às do torque de controle,
% com valores finais de $0{,}05$ Nm (eixo $x$), $0{,}47$ Nm (eixo $y$) e $-0{,}02$ Nm (eixo $z$).
% Assim, o controlador de atitude mantém a estabilização e absorve as perturbações associadas ao regime acoplado, ainda que com transitórios concentrados em instantes específicos da missão.

% \subsection{Atuação do manipulador e aproximação geométrica da ferramenta}

% O manipulador robótico opera durante a aproximação final, iniciando o modo operacional em $t=466{,}26$ s e executando a manobra até $t=900{,}00$ s.
% As posições articulares convergem para uma configuração final bem definida, com $j_1=1{,}55$ rad, $j_2=1{,}22$ rad e $j_3=0{,}23$ rad.
% As velocidades residuais ao final são $\dot{j}_1=-0{,}12$ rad/s, $\dot{j}_2=0{,}01$ rad/s e $\dot{j}_3=-0{,}19$ rad/s.

% Os torques de controle nas juntas permanecem em níveis reduzidos, na faixa de décimos de N$\cdot$m.
% Ao final, obtém-se $\tau_1=0{,}09$ N$\cdot$m, $\tau_2=8{,}95\times 10^{-4}$ N$\cdot$m e $\tau_3=0{,}04$ N$\cdot$m,
% com a junta 2 apresentando esforço residual significativamente menor em relação às demais.

% Do ponto de vista geométrico, em relação à referência de cinemática inversa para a ponta da ferramenta,
% o erro cartesiano de posição no referencial $\{0\}$ converge para valores finais de
% $e_x=1{,}00\times 10^{-2}$ m, $e_y=-4{,}95\times 10^{-3}$ m e $e_z=-1{,}04\times 10^{-2}$ m.
% A magnitude do erro de posição do efetuador final também converge, atingindo $1{,}53\times 10^{-2}$ m no instante final.

% Por fim, o alinhamento direcional entre o eixo da ferramenta e a normal da porta foi avaliado por duas métricas: (i) o ângulo entre os vetores e (ii) o produto escalar entre as direções normalizadas, ambos no referencial $\{0\}$.
% Até $t=466{,}26$ s, o ângulo permanece em $90{,}00^\circ$, indicando condição perpendicular anterior ao modo operacional do manipulador.
% Após o início do modo operacional, o ângulo é reduzido, atingindo $5{,}16^\circ$ no instante de acoplamento ($t=662{,}57$ s) e $6{,}77^\circ$ ao final.
% De forma consistente, o produto escalar evolui para valores próximos de $1$, com $0{,}9959$ no acoplamento e $0{,}9930$ ao final.
% A análise detalhada dessas métricas e de suas componentes está apresentada na Subseção~\ref{subsec:cap4_alinhamento}.


% \subsection{Síntese do alinhamento ferramenta--porta}
% Além da convergência translacional, o alinhamento direcional entre o eixo da ferramenta e a normal da porta foi avaliado por duas métricas:
% (i) o ângulo entre os vetores e (ii) o produto escalar entre as direções normalizadas, ambos em $\{0\}$.
% Até $t=466{,}26$ s, o ângulo permanece em $90{,}00^\circ$, indicando condição perpendicular anterior ao modo operacional do manipulador.
% Após o início do modo operacional, o ângulo é reduzido, atingindo $5{,}16^\circ$ no instante de acoplamento ($t=662{,}57$ s) e $6{,}77^\circ$ ao final.
% De forma consistente, o produto escalar evolui para valores próximos de $1$, com $0{,}9959$ no acoplamento e $0{,}9930$ no instante final.
% Esses resultados indicam alinhamento direcional elevado durante a fase de aproximação final e compatível com o objetivo geométrico de engate,
% ainda que a simulação não inclua modelagem detalhada da interface mecânica de captura e contato.

% \subsection{Encerramento e limitações do modelo}

% Os resultados indicam que a arquitetura proposta para a missão de \gls{rvdb} atinge os objetivos estabelecidos.
% O controle de translação baseado nas equações de \gls{hcw} conduz o perseguidor à vizinhança do alvo, e o controle de atitude mantém o alinhamento com o alvo mesmo quando o \gls{mr} está no modo operacional.
% Além disso, o manipulador robótico é capaz de posicionar a ferramenta em uma vizinhança próxima da posição da porta de acoplamento.

% Por outro lado, o modelo apresenta limitações que devem ser consideradas na interpretação dos resultados.
% A dinâmica relativa é descrita por equações linearizadas \gls{hcw} em órbitas circulares coplanares, sem qualquer consideração de perturbações orbitais, atrasos de atuadores, ruídos de sensores ou saturações detalhadas de torque.
% O manipulador é representado por um modelo rígido de três graus de liberdade, sem flexibilidade estrutural nem dinâmica de contato com a porta de acoplamento.
% Além disso, a condição de alinhamento angular entre ferramenta e porta não possui modelagem da interface mecânica de engate.